\scanpage{12}
\mspage{\unsure{6}}
Programme du \deleted{paradigme \unsure{de théorie des|de la théorie des}} \inserted{``gruppentheoretisches} Gleichnis'' (paradigme en termes de théorie des groupes) de l'\inserted{arithmétique :}\deleted{\unsure{arithm.} géométrie algébrique absolue (en car. 0).}

a) Donner une description ``gruppentheoretisch'' complète de $\Gamma$

b) Donner une description complète de tous les éléments de structure énumérés ci-dessus, dans des cas typiques importants pour $(\mathfrak X,D,J)$, dont nous indiquons quelques-uns plus bas.

\margin{On tâchera de prendre $\mathfrak X$ propre autant que possible (quitte à compléter\ldots)}

\uline{NB} Quand on a les éléments de structure envisagés pour $(\mathfrak X,D,J)$, on les a aussi automatiquement pour \unsure{tt|$\forall$} rev. \inserted{étale} fini de $\mathfrak X'....$ Il faudrait expliciter ce point avec soin, et expliciter également dans les $=$ optiques les foncteurs ``bouchage de trous'' (en rajoutant \inserted{certains} des $D_i$ qu'on avait enlevés) ou ``forage des trous'' \inserted{si $\dim\mathfrak X=1$} (en enlevant certains des $\xi\in J$).
\mspage{11}
